\bookchapter{第十章：移动支付与生活服务}

\chapterintro{说实话，我爸妈那一辈人对“手机付钱”这件事，一开始是抗拒的。“钱在手机里，被人偷了怎么办？”这个担心很正常，但现实是：你裤兜里揣着的现金，丢的概率比手机支付被盗刷的概率高得多。手机支付不是把钱存进手机，它只是一把钥匙，一把打开你银行账户的钥匙。}

\booksection{一、到底怎么绑定银行卡？}

不管你是用支付宝还是微信，第一次绑卡都要走这几步：

\begin{enumerate}
  \item 打开App，找到"我的"或"钱包"。点"添加银行卡"。
  \item 输入卡号。现在大部分App都支持直接拍照识别卡号，不用一个个数字敲。
  \item 银行会发一条验证码短信到你预留的手机号上。填进去，设置一个6位数字的支付密码。
  \item 这个支付密码和你的银行卡密码、锁屏密码，最好是三个不同的密码，别图省事设成一样的。
\end{enumerate}

有个小窍门：如果担心安全问题，可以去银行App里开一张II类或III类电子账户，专门用来绑定手机支付。这张卡里只放平时零花的钱，就算出了事，损失也有限。

\booksection{二、付款码是“谁看谁收钱”}

很多人搞不清楚付款码和收款码的区别，这个区别很重要。

\begin{itemize}
  \item \textbf{付款码}：你展示给收银员扫的那个条形码+二维码。这个码相当于你的"电子钱包"，谁扫了谁就能从你账上划走钱。所以\textbf{千万别截图发给任何人}。这个码每分钟自动刷新一次，就是为了防止被盗用。
  \item \textbf{收款码}：别人扫了这个码给你转账。这个码可以截图保存，可以发给别人，安全。
\end{itemize}

认准这个区别，能避开很多诈骗陷阱。凡是有人让你“把付款码截图发过来”的，百分之百是骗子。

\booksection{三、这些安全设置，花5分钟设好能省很多心}

\begin{itemize}
  \item \textbf{指纹/面容支付}：在支付宝和微信的支付设置里，把生物支付打开。这样每次付款要按指纹或扫脸，就算手机丢了，捡到的人也没法花你的钱。
  \item \textbf{安全锁}：支付宝和微信都有"安全锁"功能。打开之后，进入App的"钱包"页面需要额外验证一次身份。路径：支付宝 我的 设置 账号与安全 安全中心 安全锁。
  \item \textbf{免密支付检查}：很多人签约了各种免密支付，自己都忘了。去支付宝和微信的"支付设置"里，找到"免密支付/自动扣款"，看看签约了哪些服务，把不需要的关掉。
  \item \textbf{登录设备管理}：定期看看哪些设备登录过你的账号。如果发现不认识的设备，立刻踢掉。路径：支付宝 设置 账号与安全 登录设备管理。
\end{itemize}

\booksection{四、公交地铁：手机就是交通卡}

现在出门坐公交地铁，已经很少人用实体的交通卡了。手机刷闸机，有两条路可以走：

\begin{itemize}
  \item \textbf{手机钱包里的虚拟交通卡}：在苹果钱包、华为钱包、小米钱包里直接开通一张虚拟交通卡。这个方案的好处是：刷的时候不用解锁手机，直接靠上去就行，响应速度最快，而且没网也能刷。部分城市的虚拟交通卡还有和实体卡一样的换乘优惠。
  \item \textbf{支付宝/微信的乘车码}：在App里搜索"乘车码"，开通你所在城市的公交或地铁乘车码。这个方案的好处是：不用充值，直接从余额或绑定的银行卡里扣。但需要联网，而且刷之前要打开App找到码。
\end{itemize}

我个人更推荐用手机钱包的虚拟交通卡。虽然第一次要充个10块20块进去，但用起来真的方便太多，尤其是赶地铁的时候，不用停下来解锁手机找App。

\booksection{五、生活缴费、外卖、打车}

\begin{itemize}
  \item \textbf{水电气缴费}：在支付宝或微信的"生活缴费"里，把你家的水、电、燃气户号绑定好。以后每个月收到账单，直接点进去就能交。户号一般在你家的纸质账单上能找到，或者打客服电话问。
  \item \textbf{外卖比价}：同一家店在美团和饿了么上的价格和优惠经常不一样。下单前两边都看一眼，有时候能省好几块。另外，美团那个"拼好饭"功能，和朋友一起点，价格能便宜不少。
  \item \textbf{打车安全}：在滴滴或者高德打车里，记得把"紧急联系人"和"行程分享"设置好。打车的时候，App会自动把你的行程分享给你设置的联系人。另外，深夜打车优先选"出租车"而不是"快车"，出租车有行业监管，安全系数更高一些。
\end{itemize}

\booksection{六、用手机搞定医保、社保、政务}

这几年的变化很大，以前要跑大厅排队的很多事，现在手机上几分钟就办完了。

\begin{itemize}
  \item \textbf{医保电子凭证}：在支付宝或微信里搜"医保电子凭证"，激活之后，去医院看病、药店买药直接刷手机就行，不用带实体医保卡。第一次需要人脸识别验证，之后每次用的时候会生成一个动态码。
  \item \textbf{社保查询}：下载"掌上12333"App或者用支付宝里的"市民中心"，可以查社保缴纳记录、养老金账户余额、失业金申领等。
  \item \textbf{个税}：每年3月到6月的个税汇算清缴，在"个人所得税"App上就能搞定。记得看看有没有可以申报的专项附加扣除（子女教育、房贷利息、赡养老人等），这些都能帮你少交税或者退税。
\end{itemize}

\begin{tcbwarn}{记住这几条，防骗保平安}
\begin{itemize}
  \item 公检法、银行、客服，不会让你转账到"安全账户"。
  \item 任何人让你把付款码截图发过去，都是骗子。
  \item 陌生链接别点，尤其是"ETC过期""社保卡停用"这类短信。
  \item 网购退款会原路返回，不需要你提供银行卡号。
\end{itemize}

\end{tcbwarn}

\begin{tcbpractice}{现在就做：给你的支付App做一次安全体检}
\begin{enumerate}
  \item 打开支付宝和微信，确认已开启指纹/面容支付。
  \item 找到"安全锁"功能并打开。
  \item 检查"免密支付/自动扣款"列表，关掉不再需要的服务。
  \item 查看"登录设备"列表，移除不认识的设备。
  \item 激活你的医保电子凭证，下次去医院就不用带实体卡了。
\end{enumerate}

\end{tcbpractice}
